\section{Concrete Examples and Illustrations}
\label{sec:examples}

To ground the abstract formalism, we provide concrete examples of \textit{Ultimate Observables} and their associated atomicity scores, as well as illustrations of qualitative dimensions and the relational density field of Layer~2.

\subsection{Atomicity Scores for Diverse Observables}

Table~\ref{tab:examples} shows a diverse set of observables, their inferred types, and the resulting scores from the primary decomposition frameworks. The information-theoretic score for ultra-short strings (marked with an asterisk) is artificially low due to zlib compression overhead; this is a known limitation of the proxy measure, not a theoretical falsification. For longer, simple patterns, the score correctly approaches 1.0.

\begin{table}[htbp]
    \centering
    \begin{tabular}{l l c c c c}
        \toprule
        \textbf{Observable} & \textbf{Type} & \textbf{Boolean} & \textbf{Measure} & \textbf{Categorical} & \textbf{Info} \\
        \midrule
        Integer $1$ & DISCRETE & 1.00 & 1.00 & 1.00 & 0.00* \\
        Integer $42$ & DISCRETE & 1.00 & 1.00 & 1.00 & 0.00* \\
        Set $\{1,2,3\}$ & DISCRETE & 0.00 & 0.00 & 1.00 & 0.00* \\
        String ``aaaaa...'' (200x) & RELATIONAL & 1.00 & 1.00 & 1.00 & 1.00 \\
        Vector $[1.0, 0.0, 0.0]$ & CONTINUOUS & 0.50 & 1.00 & 1.00 & 0.85 \\
        Matrix $I_3$ (3x3 identity) & TOPOLOGICAL & 0.50 & 1.00 & 1.00 & 0.72 \\
        Complex $3+4i$ & QUANTUM & 1.00 & 1.00 & 1.00 & 0.00* \\
        \bottomrule
    \end{tabular}
    \caption{Examples of \textit{Ultimate Observables} and their atomicity scores. (*) The information-theoretic score for ultra-short strings is artificially low due to zlib compression overhead; this is a known limitation of the proxy measure, not a theoretical falsification. For longer, simple patterns, the score correctly approaches 1.0.}
    \label{tab:examples}
\end{table}

\subsection{Qualitative Dimensions in Practice}

The \Apeiron{} Framework supports \textit{qualitative dimensions}---intrinsic properties such as intensity, density, colour, and texture that enrich the irreducible units. Figure~\ref{fig:qualitative} illustrates how a single observable can carry multiple qualitative dimensions, each with its own atomicity test.

\begin{figure}[htbp]
    \centering
    \begin{tikzpicture}[
        dim/.style={rectangle, draw=black, fill=orange!10, minimum width=2.5cm, minimum height=0.7cm, align=center},
        arrow/.style={-{Stealth}, thick},
    ]
    \node[dim, fill=blue!15, minimum width=3.5cm] (obs) at (0,0) {Ultimate Observable\\\texttt{id="pixel\_42"}};
    
    \node[dim] (intensity) at (-3,2) {IntensityDimension\\value=0.8};
    \node[dim] (colour) at (3,2) {ColourDimension\\RGB=(1.0, 0.2, 0.2)};
    \node[dim] (texture) at (0,3.5) {TextureDimension\\GLCM contrast=0.4};
    
    \draw[arrow] (obs) -- (intensity);
    \draw[arrow] (obs) -- (colour);
    \draw[arrow] (obs) -- (texture);
    
    \node[draw, dashed, fit={(intensity) (colour) (texture)}, inner sep=0.5cm, label=below:{\small Atomicity per dimension}] {};
    
    \end{tikzpicture}
    \caption{An \textit{Ultimate Observable} enriched with qualitative dimensions. Each dimension is itself subject to atomicity testing via the \texttt{QualitativeDecompositionOperator}.}
    \label{fig:qualitative}
\end{figure}

\subsection{Relational Density Field}

Layer~2 introduces a \textit{DensityField} that captures the relational influence among observables. Figure~\ref{fig:density} shows a small network of observables with their mutual influences, illustrating how perspectives can shift based on relational density.

\begin{figure}[htbp]
    \centering
    \begin{tikzpicture}[
        node/.style={circle, draw=black, fill=blue!10, minimum size=1cm},
        edge/.style={-{Stealth}, thick, color=gray},
    ]
    \node[node] (a) at (0,0) {A};
    \node[node] (b) at (3,0) {B};
    \node[node] (c) at (1.5,2) {C};
    \node[node] (d) at (1.5,-1.5) {D};
    
    \draw[edge] (a) -- (b) node[midway, above] {0.8};
    \draw[edge] (a) -- (c) node[midway, left] {0.3};
    \draw[edge] (b) -- (c) node[midway, right] {0.6};
    \draw[edge] (c) -- (d) node[midway, right] {0.9};
    \draw[edge] (a) -- (d) node[midway, below] {0.2};
    
    \node[draw, dashed, fit={(a) (b) (c)}, inner sep=0.7cm, label=above:{\small Perspective A dominant}] {};
    
    \end{tikzpicture}
    \caption{A relational density field among four observables. Edge weights represent influence strength. The dashed region indicates a cluster where perspective A becomes dominant due to cumulative influence exceeding a threshold.}
    \label{fig:density}
\end{figure}

The \texttt{DensityField} class implements a two-channel influence model (resonance + embedding) with distance-based decay. The adjacency function $A(o_i, o_j) \to [0,1]$ is computed as the cosine similarity of relational embeddings. When the cumulative influence of a particular observer perspective exceeds the configured threshold, the target observable's perspective shifts, demonstrating the emergent, self-organizing nature of the relational layer.